\documentclass[11pt,a4paper]{article}
\usepackage[margin=24mm]{geometry}
\usepackage{amsmath,amssymb,amsthm,booktabs}
\newtheorem{lemma}{Lemma}\newtheorem{theorem}{Theorem}
\title{Three counts with an element of order twelve}
\author{Human Compression project}\date{}
\begin{document}\maketitle
Put $G=O'N$, $g=|G|=460815505920$, and let $x_J(p,q)$ count ordered
pairs $(r,s)\in J^2$ with $|r|=p$, $|s|=q$, $|rs|=2$. Write
$e_J(p,q)$ for the number which generate $J$. Throughout this paper
$p\in\{4,6,12\}$ and $q=12$. We prove the fixed Connor--Leemans
Table~1 entries
\[
 e_G(4,12)/(2g)=120,\quad e_G(6,12)/(2g)=474,\quad
 e_G(12,12)/(2g)=1072.                                      \tag{1}
\]
The complete subgroup roster, the primary ATLAS group structures and
centralizer orders, and the acquired character evaluations are borrowed
premises. Row numbers refer to that roster. Let $n_i$ be the number of
conjugates in row $i$. An identity $n_i k=n_j$ for a listed cover of
multiplicity $k$ proves that every actual row $j$ subgroup has a unique
row $i$ overgroup. We compose such proved uniqueness statements; we do
not multiply arbitrary cover paths.

Use the following names for maximal subgroups:
\[
\begin{array}{c|c|r}
\text{symbol}&\text{row}&\text{order}\\\hline
M,M'&2,3&3753792\\ A&7&161280\\
B&11&25920\\ C&12&25920\\ E&21&10752.
\end{array}                                                \tag{2}
\]
Here $M,M'\cong L_3(7):2$, $A$ is an involution centralizer,
$B=(F\times A_6).2$, with $F=3^2:C_4$, and
$C=3^4:(2^{1+4}.D_{10})$. The remaining maximal classes have zero
character count for all three types. The accepted order twelve component
identifies $E=4^3.L_3(2)$ with its retained literal model.

\begin{lemma}\label{distinguished}
Every element $s$ of order twelve lies in exactly one maximal subgroup
of type $A$, namely $A(s)=C_G(s^6)$. A subgroup containing $s$ belongs
to the descendant list of row 7 if and only if it lies in $A(s)$.
\end{lemma}
\begin{proof}
The acquired centralizer $C_G(s)$ has order 36. Its central subgroup
$\langle s^3\rangle$ is a Sylow 2-subgroup, and hence the unique one.
It has a unique involution $s^6$. Thus the only involution centralizer
containing $s$ is $C_G(s^6)$. The last assertion follows by conjugating
an actual containment represented in the complete roster.
\end{proof}
We call a pair covered by $A$ when it lies in this unique $A(s)$.
Within any group $J$, its covered pairs are precisely those for which
$r$ commutes with $s^6$.

\begin{lemma}\label{B}
Put $B_0=F\times A_6$, the row 17 subgroup of index two in $B$.
Every pair under consideration in $B$ lies in $B_0$. The numbers of
pairs in one $B$ which are covered by $A$ are
\[
 b_A(p)=9x_{C_4\times A_6}(p,12)+45x_{F\times D_8}(p,12).
                                                               \tag{3}
\]
\end{lemma}
\begin{proof}
Let $V=O_3(F)$ and $N=A_6$. These are characteristic factors of $B_0$:
$N$ is its final derived subgroup and $V$ is its largest normal
3-subgroup. They are therefore normal in $B$. The quotient $B/(VN)$
has order eight. The roster lists exactly three subgroups of index two
in $B$, rows 17, 18 and 19, each with just one subgroup of index two,
row 25. Since $V$ has odd order and $N$ is perfect, both lie in every
such subgroup. Each index two subgroup of the quotient is consequently
cyclic of order four. Thus $B/(VN)\cong Q_8$.

The action of $C_4$ on $V$ is fixed point free away from the identity.
Indeed, for $1\ne v\in V$, the subgroup $V\times N$ already has the
acquired full centralizer order $3240$. No element outside it can fix
$v$. In particular the involution of $C_4$ acts as $-I$. Its generator
has minimal polynomial $X^2+1$ over $\mathbb F_3$, so this is the
explicit affine group $F$ used in the certificate. Its element counts
in orders $1,2,3,4$ are respectively $1,9,8,18$.

In $A_6$ there are 80 elements of order three and 90 of order four.
Hence $B_0$ has $18\cdot80+8\cdot90=2160$ elements of order twelve.
This equals the acquired total for $B$, so all of them lie in $B_0$.
The image of an involution in $Q_8$ is central. Since the image of
$s\in B_0$ belongs to a cyclic subgroup of order four, so does the
image of $r=(rs)s^{-1}$. Thus $r\in B_0$ as well.

The order twelve elements in $F\times A_6$ have component orders
$(4,3)$ or $(3,4)$. Their sixth powers are respectively one of the nine
involutions in $F$, or one of the 45 involutions in $A_6$. The relevant
centralizers in $B_0$ are $C_4\times A_6$ and $F\times D_8$.
Summing their pair counts proves (3); the unique sixth power prevents
any repeated pair.
\end{proof}

\begin{lemma}\label{affine}
In one $C$, the pairs covered by $A$ number
\[
                  c_A(p)=90x_{F\times C_4}(p,12).              \tag{4}
\]
Every pair in $C$ which is not covered by $A$ has a unique maximal
subgroup of type $C$ containing it.
\end{lemma}
\begin{proof}
Write $C=W\rtimes K$, where $W\cong\mathbb F_3^4$ and $K$ contains
an extraspecial normal subgroup $Q$ of order 32, with $K/Q\cong D_{10}$.
We reuse the accepted affine action argument: an element of order five
acts irreducibly on $W$, and the square form on $Q/Z(Q)$ has five
nonzero singular vectors. Here are the additional involution details.
The central involution $z$ of $Q$ acts as $-I$ on $W$. It commutes with
the irreducible order five action, so its action lies in the field
centralizer and is $I$ or $-I$. If it were $I$, the action of $Q$ would
factor through the elementary abelian group $Q/Z(Q)$ and split into
characters over $\mathbb F_3$. The order five action permutes its
nontrivial characters in orbits of length five; there are only four
dimensions. Thus $Q$ would act trivially. But then $Q\leq C_G(w)$ for
any $0\ne w\in W$, impossible since 32 does not divide 3240.

Each of the ten noncentral involutions $t\in Q$ is conjugate in $Q$
to $zt$. Since $z$ acts as $-I$, the positive and negative eigenspaces
of $t$ on $W$ both have dimension two. Its affine lifts which are
involutions number nine. The five singular vectors are transitive
under the order five action, and $Q$ exchanges the two lifts of each.
Consequently these 90 affine involutions form one $C$-class. Their
centralizers have order $25920/90=288$.

If $|s|=12$, its image in $K/Q\cong D_{10}$ has order at most two,
so the image of $s^6$ lies in $Q$. It cannot be $z$: the lift of $z$
has no nonzero fixed vector in $W$, whereas $s^4$ is a nonidentity
element of $W$ of order three. Thus $s^6$ belongs to that class of 90
involutions. Its centralizer in $C$ contains $s$. The complete common
ancestor and descendant calculation for a cyclic order twelve subgroup,
rows 7 and 12 has only row 197 as a maximal class. Since the actual
centralizer has order 288, it is row 197.

A row 197 subgroup lies in a row 32 subgroup, and hence in a $B$.
Its intersection with $B_0$ has order 144: a row 197 subgroup is
not below row 57, the intersection of that row 32 subgroup with $B_0$.
It contains all the order twelve pairs, by Lemma~\ref{B}. The roster identifies this intersection
as row 254, inside row 57, which is $A_6\times C_4$. Projection of an
order 144 subgroup to $A_6$ has order at least 36; there is no subgroup
of order 72 in $A_6$. It therefore has image of order 36 and kernel
$C_4$, and equals $F\times C_4$. This proves (4).

For uniqueness, the complete generated-subgroup list in $C$, outside
row 7, is
\[
 12,26,36,38,40,69,70,71,120,178,215,288.                      \tag{5}
\]
The following actual parents are unique by direct incidence counts:
\[
\begin{gathered}
26\to12,\quad36,38,40\to26,\quad69,70,71\to36,\\
120\to69,\quad178,215\to120,\quad288\to178.
\end{gathered}                                               \tag{6}
\]
For example $n_{36}\,2=n_{70}$ and $n_{120}\,12=n_{215}$.
Composing (6) proves the claim for every row in (5).
\end{proof}

\begin{lemma}\label{BC}
Let $D$ denote a row 69 subgroup. Then $D\cong F\times F$.
Pairs contained in both a $B$ and a $C$, but not covered by $A$,
lie in a unique actual $D$. Each such pair has exactly two $B$
overgroups and one $C$ overgroup. The number of these pairs in one $D$
is
\[
       d(p)=x_{F\times F}(p,12)-18x_{F\times C_4}(p,12).       \tag{7}
\]
\end{lemma}
\begin{proof}
A maximal subgroup of the direct product $F\times A_6$ is a product:
there is no common nontrivial simple quotient of its solvable and
simple factors. Its maximal subgroups have the structures
\[
\begin{array}{c|ccccc}
\text{row}&25&47&57&69&89\\\hline
\text{structure}&(3^2:C_2)\times A_6&F\times A_5&C_4\times A_6&
F\times F&F\times S_4.
\end{array}                                                  \tag{8}
\]
The two $A_5$ classes and the two $S_4$ classes in $A_6$ are combined
in the indicated global rows. In particular (8) identifies $D$.
It has exactly two $B_0$ parents, since $n_{17}\,10=2n_{69}$.
The elements of order twelve generate $F\times F$: comparing
$(a,b)$ and $(a,b^{-1})$, with orders four and three, gives each
normal $3^2$ factor, and then both cyclic order four quotients.
By Lemma~\ref{B}, any $B$ containing $D$ contains it in $B_0$.
Each $B_0$ has a unique $B$ parent, so $D$ has exactly two $B$ parents.
Its unique $C$ parent follows from (6).

Outside $A$, the complete common list for $B$ and $C$ is
\[
           36,69,70,71,120,178,215,288.                       \tag{9}
\]
Rows 36, 70 and 71 are not descendants of row 17, so cannot be
generated by our pairs in $B$. Every one of the remaining rows has
a unique global row 69 overgroup, by the incidences in (6).
We must also check that this overgroup lies in the particular $B$;
a class containment alone would not prove this.

Let $L$ be one of these actual generated subgroups in $B_0$ and
project to $A_6$. For orders 1296, 648 and 324, its image has order,
respectively, 36; 18 or 36; or 9, 18 or 36. These follow from the fact that the kernel order divides 36 and
the complete subgroup list below the roster's row 163 copy of $A_6$.
Its maximal orders are $60,60,36,24,24$, with a single class of order 36.
Each such image lies in a unique subgroup of order 36, the stabilizer
of an unordered partition of six letters into two sets of three.
For order nine this is its normalizer; for order eighteen its Sylow
3-subgroup is normal; the order 36 case is immediate.

The row 215 group is $F\times S_3$: it occurs as a maximal subgroup
of $F\times S_4$, and the roster has just one maximal row of that
order there. Its derived subgroup is a 3-group and its Sylow 2-subgroup
is abelian. Its image in $A_6$ has order at least six, and cannot be
$A_4$ or $S_4$, by these two properties. The remaining possibilities
are $S_3$ and the orders 18 and 36 just handled. Every $S_3$ in
$A_6$ determines a unique invariant partition into two triples:
its orbit lengths are $3+3$ or $3+2+1$. A transitive action of degree
six would be regular and give odd permutations for its involutions.
Similarly row 288 is $F\times C_3$, as realized inside $F\times F$;
it is the only row of order 108 containing an element of order twelve.
Its derived subgroup is a 3-group, excluding an $A_4$ image. The
possible image orders are 3, 6, 9, 18 and 36. A subgroup of order
three also determines a unique partition into triples from its cycles.
Thus every such $L$ lies in a row 69 subgroup inside this $B_0$.
Global uniqueness forces it to be the same $D$ as before. Since $D$
has a unique $C$ parent, that parent must also be the original $C$.
This proves all the actual containment and multiplicity assertions.

Finally, in $F\times F$, an element of order twelve has component
orders three and four. Its sixth power is an involution in exactly
one factor. There are nine choices in either factor, each with
centralizer $F\times C_4$. This proves (7).
\end{proof}

\begin{lemma}\label{Bunique}
A pair in $B$ which is outside all $A$ and $C$ subgroups has a unique
$B$ overgroup.
\end{lemma}
\begin{proof}
After the exclusions in Lemma~\ref{B}, the possible generated rows
outside $A$ and $C$ reduce to
\[
        17,25,31,47,48,87,89,146,147,214,251,253,319,320.
\]
Rows 17, 25, 31, 48 and 87 have a quotient $A_6$. For rows 31, 48
and 87, this also follows from their respective indices two, three
and three in rows 25, 25 and 31: the perfect simple factor $A_6$
acts trivially on such a small coset set and is contained in the subgroup.
Neither of the
projected generator orders can exceed four, and neither can be five.
A group generated by orders at most four whose product has order at
most two is, in these cases, a quotient of a dihedral group or of
one of the triangle groups $(2,3,3)$, $(2,3,4)$, $(2,4,4)$.
These groups are solvable, so cannot map onto $A_6$. Row 47 has a
quotient $A_5$, whose permitted projected orders here are at most three;
the same argument excludes it. This reuses the accepted triangle
exclusion, including its elementary proof for $(2,4,4)$.

The two order 72 rows 319 and 320 are $C_3\times S_4$, as identified
in the next lemma, and their generating counts for our three types
are zero. For the remaining rows, direct incidences give unique parents
\[
 89\to17\to11,\quad146,147\to89,\quad214\to146,
 \quad251,253\to146.
\]
For instance $n_{17}\,30=n_{89}$ and $n_{146}\,6=n_{251}=n_{253}$.
Their compositions prove the claim.
\end{proof}

\begin{lemma}\label{projective}
Write $T$ for row 22, $H$ for the retained block group of order 1344,
and $T_A(p)$ for the pairs in $T$ covered by $A$. The union of pairs
in the two projective maximal families, outside $A$, has size
\[
 P(p)=2n_2\bigl(x_M(p,12)-2793x_H(p,12)\bigr)
                  -n_{22}\bigl(x_T(p,12)-T_A(p)\bigr).       \tag{10}
\]
Its intersection with the $B$ family outside $A$ has size
$2n_{251}\epsilon(p)$, where $\epsilon(4)=\epsilon(12)=0$ and
$\epsilon(6)=144$. It has no intersection there with the $C$ or $E$
families.
\end{lemma}
\begin{proof}
For each projective maximal, the accepted centralizer argument gives
2793 inner involutions with centralizer $H$. Since $s^6$ is inner,
this accounts for its $A$-covered pairs.
Outside $A$ the complete generated lists in the two maximal classes
are $2,22,29,251,319$ and $3,22,29,253,320$.
Every actual row 22 subgroup has one parent of each projective class,
by $n_2\,456=n_{22}=n_3\,456$. Row 29 has a unique row 22 parent.
Rows 251 and 253 each have a unique projective parent of the indicated
class, by $n_2\,26068=n_{251}$ and its counterpart. The order 72
rows have zero generating counts. Thus exactly the pairs in $T$
outside $A$ have been counted twice, proving (10).

Here is a structural identification of the order 144 rows used in
this argument. In $L_3(7)$ the diagonal torus is $C_3\times C_2^2$.
The Weyl group $S_3$ acts on the second factor in its usual way, and
acts on $C_3$ by its sign character. Inverse transpose inverts $C_3$
and fixes the $C_2^2$ and permutation factors. Twisting the $S_4$
complement by this outer involution according to its sign identifies
the full monomial group with $S_3\times S_4$.
This group contains an outer element and an element of order twelve.
The complete maximal orders of $L_3(7):2$ leave only an order 144
maximal for it: the socle cannot contain the outer element, and none
of the other proper maximal orders is divisible by 144. Among the
two order 144 classes only row 251 (respectively 253) contains order
twelve. This identifies the actual row, not merely its character.

In $S_3\times S_4$, only the index two subgroup $C_3\times S_4$
contains order twelve elements. The other two have respectively no
order four component, or incompatible parities. Hence the order 72
rows containing twelve are the asserted products. Literal multiplication
and generated closures in these groups give
\[
\begin{array}{c|rrr}
 &p=4&p=6&p=12\\\hline
 e_{S_3\times S_4}(p,12)&0&144&0\\
 e_{C_3\times S_4}(p,12)&0&0&0.
\end{array}
\]
The other generated orders in the first row's full pair sets are 24
and 12, which are covered by $A$ in the roster. Each actual row 251
or 253 has a unique $B$ parent by Lemma~\ref{Bunique}. The two global
classes are disjoint, giving the stated intersection. The other
intersection assertions follow from the disjoint complete generated
lists outside $A$.
\end{proof}

\begin{lemma}\label{E}
The union of pairs in the $E$ family outside $A$ has size
\[
             E(p)=n_{21}\bigl(x_E(p,12)-7x_J(p,12)\bigr),     \tag{11}
\]
where $J$ is the retained centralizer of order 1536.
\end{lemma}
\begin{proof}
The accepted model and order twelve proof establish that every sixth
power lies in the single orbit of seven nonidentity involutions of
the normal $C_4^3$. Their centralizers are $J$, proving the subtraction.
The possible generated rows outside $A$ are $21,56,68,98$.
Row 21 has quotient $L_3(2)$, which has no elements of order six or
twelve; the permitted projected orders are again at most four and
cannot generate this nonsolvable group by the triangle exclusion.
Row 68 has the noncyclic odd quotient $7:3$. The images of the two
generators there would be inverse, since their product has order at
most two, so cannot generate it. Row 56 has a unique $E$ parent and
row 98 a unique row 56 parent, by direct incidences. This proves that
all remaining outside pairs have a unique $E$ parent, as required.
\end{proof}

\begin{theorem}
The three counts in (1) follow from the above lemmas and the exact
small-group values in the following table:
\[
\begin{array}{c|rrr}
J&p=4&p=6&p=12\\\hline
H\ (1344)&672&6496&112\\
T\ (8232)&0&268912&67228\\
T_A&0&38416&1372\\
J\ (1536)&3072&3072&7936\\
C_4\times A_6&5760&0&3040\\
F\times D_8&288&64&16\\
F\times C_4&288&0&16\\
F\times F&25920&0&2592
\end{array}                                                   \tag{12}
\]
\end{theorem}
\begin{proof}
The number of pairs in the union of the $B$ and $C$ families outside
$A$ is, by Lemmas~\ref{affine}--\ref{Bunique},
\[
 Q(p)=n_{12}(x_C-c_A)+n_{11}(x_B-b_A)-2n_{69}d(p).             \tag{13}
\]
Indeed each common outside pair is counted once in the $C$ term and
twice in the $B$ term, and every other $B$ pair is counted just once.
All $E$ intersections outside $A$ with the other families are empty
by their generated-subgroup lists. It follows that the complete
nongenerating count is
\[
 n_7x_A+P(p)+Q(p)+E(p)-2n_{251}\epsilon(p).                   \tag{14}
\]
For clarity, the five terms in (14), divided by $g$, are
\[
\begin{array}{c|rrrrr}
p&n_7x_A/g&P/g&Q/g&E/g&2n_{251}\epsilon/g\\\hline
4&7/2&56&8&0&0\\
6&46/9&144&2&4&2\\
12&73/36&64&40/9&0&0.
\end{array}
\]
Subtract (14) from the acquired $x_G$ and divide by
$|\operatorname{Aut}(G)|=2g$. The action on generating pairs is free,
so the results are exactly (1), with no further duality division.

The first four lines of (12) reuse already acquired complete literal
models and their complete order twelve class and partner lists. The
remaining lines use explicit affine permutations of $\mathbb F_3^2$,
alternating permutations on six letters and cyclic or dihedral
permutations on four letters. The independent reader checks complete
group closures, all order twelve classes and all involution partners;
it also checks the generated closures used in Lemma~\ref{projective}.
All numerical operations and source comparisons are exact.
\end{proof}

The numerical assertions are due to Connor and Leemans. The contribution
here is the common reproof by actual subgroup incidences, the restriction
to $B_0$ and the explicit small product counts. The restriction argument
was developed jointly with the parent agent; the final review is distinct.
No new chirality statement, whole-source completion, open-conjecture
solution or Lean verification is claimed by this component.
\end{document}
